% Unit 041 English translation of chapter3.tex lines 1710--1881.
% Source: Methods of Algebra, Volume 2, commit
% 9a5803ff77dd3257484cb177f851a73770a59dd3 (CC BY 4.0).
% stable-unit-id: o014.aljabr2.chapter3.double-complex-cohomology
% Coverage: Section 3.10 in full.

% segment-id: o014.li2.chapter3.double-complex-cohomology.h001
\section{Cohomology of double complexes}\label{sec:double-cplx-coh}
\hypertarget{o014.aljabr2.chapter3.double-complex-cohomology}{ }

% segment-id: o014.li2.chapter3.double-complex-cohomology.p001
This section examines the cohomology of a double complex $X$ in the horizontal
and vertical directions, and its relationship with the cohomology of the total
complex. The results will be used to study derived functors of bifunctors; the
argument follows \cite[\S 12.5]{KS06}.

% segment-id: o014.li2.chapter3.double-complex-cohomology.p002
Let $\mathcal{A}$ be an abelian category. As always, when $\mathcal{A}$ is a
$\Bbbk$-linear abelian category, all the results have corresponding
generalizations.

% segment-id: o014.li2.chapter3.double-complex-cohomology.p003
Recall that Definition~\ref{def:bicplx-F1} (or the more explicit diagram
\eqref{eqn:bicplx-diagram}) introduced a pair of invertible additive functors
% segment-id: o014.li2.chapter3.double-complex-cohomology.g001
\[\begin{tikzcd}
	\cate{C}^2(\mathcal{A}) \arrow[r, shift left, "{F_{\mathrm{I}}}"] \arrow[r, shift right, "{F_{\mathrm{II}}}"'] & \cate{C}(\cate{C}(\mathcal{A})) .
\end{tikzcd}\]
Using them, define for every double complex $X$ (with $p,q,n\in\Z$ below)
% segment-id: o014.li2.chapter3.double-complex-cohomology.q001
\begin{align*}
	\Hm_{\mathrm{I}}^p(X) & := (\Hm^p \circ F_{\mathrm{I}})(X)
	= \text{the complex}\; \left[ \begin{tikzcd}[row sep=small]
		\vdots \\
		\Hm^p(X^{\bullet, q+1}, \dhori) \arrow[u, "{\dvert}" inner sep=0.6em] \\
		\Hm^p(X^{\bullet, q}, \dhori) \arrow[u, "{\dvert}" inner sep=0.6em] \\
		\vdots \arrow[u, "{\dvert}" inner sep=0.6em]
	\end{tikzcd}\right] \quad \text{(see Proposition~\ref{prop:abelian-cat-cplx})}, \\
	\Hm_{\mathrm{II}}^q(X) & := (\Hm^q \circ F_{\mathrm{II}})(X) \\
	& = \text{the complex}\; \left[ \cdots \xrightarrow{\dhori} \Hm^q(X^{p, \bullet}, \dvert) \xrightarrow{\dhori} \Hm^q(X^{p+1, \bullet}, \dvert) \xrightarrow{\dhori} \cdots \right] , \\
	\tau_{\mathrm{I}}^{\leq n} & := (F_{\mathrm{I}})^{-1} \circ \tau^{\leq n} \circ F_{\mathrm{I}}, \\
	\tau_{\mathrm{II}}^{\leq n} & := (F_{\mathrm{II}})^{-1} \circ \tau^{\leq n} \circ F_{\mathrm{II}},
\end{align*} \index[sym1]{HmIpX@$\Hm_{\mathrm{I}}^p(X)$, $\Hm_{\mathrm{II}}^q(X)$} \index[sym1]{tauI@$\tau_{\mathrm{I}}^{\leq n}$, $\tau_{\mathrm{II}}^{\leq n}, \ldots$}%
\footnote{Translator's note (O014-C054): in the first two definitions the
source prints only the composite functors $\Hm^p\circ F_{\mathrm{I}}$ and
$\Hm^q\circ F_{\mathrm{II}}$ on the right, although the left-hand sides are
their values on the double complex $X$. The argument $X$ has been restored.}
There are likewise functors $\tilde{\tau}_{\mathrm{I}}^{\leq n}$,
$\tau_{\mathrm{I}}^{\geq n}$, $\tilde{\tau}_{\mathrm{I}}^{\geq n}$ and
$\tilde{\tau}_{\mathrm{II}}^{\leq n}$,
$\tau_{\mathrm{II}}^{\geq n}$, $\tilde{\tau}_{\mathrm{II}}^{\geq n}$.
Thus there are still morphisms of functors
% segment-id: o014.li2.chapter3.double-complex-cohomology.q002
\begin{gather*}
	\tau_{\mathrm{I}}^{\leq n} \to \tilde{\tau}_{\mathrm{I}}^{\leq n} \to \identity_{\cate{C}^2(\mathcal{A})} \to \tilde{\tau}_{\mathrm{I}}^{\geq n} \to \tau_{\mathrm{I}}^{\geq n}, \\
	\tau_{\mathrm{II}}^{\leq n} \to \tilde{\tau}_{\mathrm{II}}^{\leq n} \to \identity_{\cate{C}^2(\mathcal{A})} \to \tilde{\tau}_{\mathrm{II}}^{\geq n} \to \tau_{\mathrm{II}}^{\geq n}.
\end{gather*}

% segment-id: o014.li2.chapter3.double-complex-cohomology.de001
\begin{definition}\label{def:Cf-Supp}
	\index[sym1]{SuppX@$\Supp(X)$}
	\index[sym1]{C2f@$\cate{C}^2_f(\mathcal{A})$}
	For a double complex $X$, write
	$\Supp(X) := \left\{ (p,q) \in \Z^2 : X^{p,q} \neq 0 \right\}$.
	Let $\cate{C}^2_f(\mathcal{A})$ be the full subcategory consisting of
	those objects $X$ of $\cate{C}^2(\mathcal{A})$ for which the set
	$\{(p,q)\in\Supp(X):p+q=n\}$ is finite for every $n\in\Z$.
\end{definition}

% segment-id: o014.li2.chapter3.double-complex-cohomology.ex001
\begin{example}
	If $\Supp(X)\subset\Z_{\geq 0}^2$ (the first quadrant), if
	$\Supp(X)\subset\Z_{\leq 0}^2$ (the third quadrant), or if $\Supp(X)$
	is the union of finitely many rows or columns, then $X$ is an object of
	$\cate{C}^2_f(\mathcal{A})$.
\end{example}

% segment-id: o014.li2.chapter3.double-complex-cohomology.p004
The total complexes $\tot_{\oplus}(X)$ and $\tot_{\Pi}(X)$ are defined for
every object $X$ of $\cate{C}^2_f(\mathcal{A})$, with no additional assumption
on $\mathcal{A}$. Moreover, $\tot_{\oplus}(X)=\tot_{\Pi}(X)$ in this case.
We therefore obtain an additive functor
$\tot:\cate{C}^2_f(\mathcal{A})\to\cate{C}(\mathcal{A})$.

% segment-id: o014.li2.chapter3.double-complex-cohomology.p005
We now make a simple but necessary observation. Given an object $X$ of
$\cate{C}^2_f(\mathcal{A})$ and $n\in\Z$, there is a finite subset
$S\subset\Z^2$ such that $\Hm^n(\tot(X))$ is completely determined by the
data $\left(X^{p,q},\dhori^{p,q},\dvert^{p,q}\right)_{(p,q)\in S}$.

% segment-id: o014.li2.chapter3.double-complex-cohomology.le001
\begin{lemma}\label{prop:tot-exact}
	The total-complex functor
	$\tot:\cate{C}^2_f(\mathcal{A})\to\cate{C}(\mathcal{A})$ is exact
	(Definition~\ref{def:exact-functor}).
\end{lemma}
% segment-id: o014.li2.chapter3.double-complex-cohomology.pf001
\begin{proof}
	Let $X\xrightarrow{f}Y\xrightarrow{g}Z$ be an exact sequence in
	$\cate{C}^2_f(\mathcal{A})$. Recalling the relevant definitions, such as
	Proposition~\ref{prop:abelian-cat-cplx}, it suffices to show that for every
	$n\in\Z$ the sequence
% segment-id: o014.li2.chapter3.double-complex-cohomology.q003
	\[ \bigoplus_{p+q=n} X^{p,q} \xrightarrow{\bigoplus_{p+q=n} f^{p,q}} \bigoplus_{p+q=n} Y^{p,q} \xrightarrow{\bigoplus_{p+q=n} g^{p,q}} \bigoplus_{p+q=n} Z^{p,q} \]
	is exact in $\mathcal{A}$. Since these direct sums are finite, this reduces
	to the exactness of
	$X^{p,q}\xrightarrow{f^{p,q}}Y^{p,q}\xrightarrow{g^{p,q}}Z^{p,q}$.
\end{proof}

% segment-id: o014.li2.chapter3.double-complex-cohomology.p006
Although the next lemma may look complicated, it is merely an exercise in the
definitions and applies to every additive category $\mathcal{A}$.

% segment-id: o014.li2.chapter3.double-complex-cohomology.le002
\begin{lemma}\label{prop:double-cplx-tot-aux0}
	Let $h\in\Z$ and let $Y$ be an object of $\cate{C}(\mathcal{A})$. Form
% segment-id: o014.li2.chapter3.double-complex-cohomology.l001
	\begin{compactitem}
% segment-id: o014.li2.chapter3.double-complex-cohomology.i001
		\item the object $Y[-h]$ of $\cate{C}(\mathcal{A})$;
% segment-id: o014.li2.chapter3.double-complex-cohomology.i002
		\item the object $Y_{\mathrm{I}}[-h]$ of
		$\cate{C}(\cate{C}(\mathcal{A}))$, whose term in degree $h$ is $Y$
		and whose other terms are $0$.
	\end{compactitem}
	There are then canonical isomorphisms in $\cate{C}(\mathcal{A})$
% segment-id: o014.li2.chapter3.double-complex-cohomology.q004
	\begin{gather*}
		(\tot \circ F_{\mathrm{I}}^{-1}) \Cone\left( \identity_{Y_{\mathrm{I}}[-h]} \right) \simeq \Cone\left( \identity_{Y[-h]} \right), \\
		(\tot \circ F_{\mathrm{I}}^{-1}) \left( Y_{\mathrm{I}}[-h] \right) \simeq Y[-h].
	\end{gather*}
\end{lemma}
% segment-id: o014.li2.chapter3.double-complex-cohomology.pf002
\begin{proof}
	By definition,
	$\Cone\left(\identity_{Y_{\mathrm{I}}[-h]}\right)$ is the object
	$Y\xrightarrow{\identity_Y}Y$ of
	$\cate{C}^{[h-1,h]}(\cate{C}(\mathcal{A}))$, concentrated in degrees
	$h-1$ and $h$. Hence
	$Z:=F_{\mathrm{I}}^{-1}\Cone\left(\identity_{Y_{\mathrm{I}}[-h]}\right)$
	is the double complex shown below:
% segment-id: o014.li2.chapter3.double-complex-cohomology.g002
	\[\begin{tikzcd}
		n \in \Z & \vdots & \vdots \\
		(q = n-h+1) & Y^{n-h+1} \arrow[u, "d_Y"] \arrow[r, "\identity"] & Y^{n-h+1} \arrow[u, "d_Y"] \\
		(q = n-h) & Y^{n-h} \arrow[r, "\identity"] \arrow[u, "d_Y"] & Y^{n-h} \arrow[u, "d_Y"] \\
		& (p = h-1) \arrow[phantom, u, "\cdots" description, sloped] & (p = h) \arrow[phantom, u, "\cdots" description, sloped]
	\end{tikzcd}\]
	All columns corresponding to $p\notin\{h-1,h\}$ are zero. The definition
	of the total complex (Definition~\ref{def:total-cplx}) now gives at once
% segment-id: o014.li2.chapter3.double-complex-cohomology.q005
	\begin{align*}
		\tot(Z)^n & = Y^{n-h+1} \oplus Y^{n-h} = Y[-h]^{n+1} \oplus Y[-h]^n \\
		& = \Cone\left( \identity_{Y[-h]} \right)^n , \\
		d_{\tot(Z)}^n & = \begin{pmatrix}
			(-1)^{h-1} d_Y^{n-h+1} & 0 \\
			\identity_{Y^{n-h+1}} & (-1)^{h} d_Y^{n-h}
		\end{pmatrix}
		= \begin{pmatrix}
			- d_{Y[-h]}^{n+1} & 0 \\
			\identity_{Y[-h]^{n+1}} & d_{Y[-h]}^n
		\end{pmatrix} \\
		& = d_{\Cone\left( \identity_{Y[-h]} \right)}^n .
	\end{align*}
	The second isomorphism can be checked similarly, but more easily, and is
	left to the reader.
\end{proof}

% segment-id: o014.li2.chapter3.double-complex-cohomology.le003
\begin{lemma}\label{prop:double-cplx-tot-aux}
	Let $X$ be an object of $\cate{C}^2_f(\mathcal{A})$ and let $q\in\Z$.
% segment-id: o014.li2.chapter3.double-complex-cohomology.l002
	\begin{enumerate}[(i)]
% segment-id: o014.li2.chapter3.double-complex-cohomology.i003
		\item The natural morphism
		$\tot\left(\tau_{\mathrm{I}}^{\leq q}X\right)\to
		\tot\left(\tilde{\tau}_{\mathrm{I}}^{\leq q}X\right)$ is a
		quasi-isomorphism.
% segment-id: o014.li2.chapter3.double-complex-cohomology.i004
		\item There is a canonical short exact sequence
% segment-id: o014.li2.chapter3.double-complex-cohomology.q006
		\[ 0 \to \tot\left( \tilde{\tau}_{\mathrm{I}}^{\leq q-1} (X) \right) \to \tot\left( \tau_{\mathrm{I}}^{\leq q}(X) \right) \to \Hm_{\mathrm{I}}^q(X)[-q] \to 0. \]
	\end{enumerate}
\end{lemma}
% segment-id: o014.li2.chapter3.double-complex-cohomology.pf003
\begin{proof}
	Apply Lemma~\ref{prop:truncation-ses} in
	$\cate{C}(\cate{C}(\mathcal{A}))$ to obtain the short exact sequences
% segment-id: o014.li2.chapter3.double-complex-cohomology.q007
	\begin{gather*}
		0 \to \tau^{\leq q} F_{\mathrm{I}} (X) \to \tilde{\tau}^{\leq q} F_{\mathrm{I}} (X) \to \Cone\left( \identity_{\Image\left( d^q_{F_{\mathrm{I}} X} \right)_{\mathrm{I}} [-q-1]} \right) \to 0, \\
		0 \to \tilde{\tau}^{\leq q-1} F_{\mathrm{I}}(X) \to \tau^{\leq q} F_{\mathrm{I}}(X) \to \Hm^q_{\mathrm{I}}(X)_{\mathrm{I}}[-q] \to 0.
	\end{gather*}
	The notation $(\cdots)_{\mathrm{I}}[\cdots]$ has the meaning specified in
	Lemma~\ref{prop:double-cplx-tot-aux0}. Apply
	$\tot\circ F_{\mathrm{I}}^{-1}$ to these short exact sequences.
	Lemma~\ref{prop:tot-exact} says that $\tot$ is exact, and
	$F_{\mathrm{I}}^{-1}$ is of course exact as well. Substitution from
	Lemma~\ref{prop:double-cplx-tot-aux0} gives short exact sequences in
	$\cate{C}(\mathcal{A})$
% segment-id: o014.li2.chapter3.double-complex-cohomology.q008
	\begin{gather*}
		0 \to \tot\left(\tau_{\mathrm{I}}^{\leq q} X\right) \to \tot\left( \tilde{\tau}_{\mathrm{I}}^{\leq q} X \right) \to \Cone\left( \identity_{\Image( d_{F_{\mathrm{I}} X}^q )[-q-1]} \right) \to 0, \\
		0 \to \tot\left( \tilde{\tau}_{\mathrm{I}}^{\leq q-1} X \right) \to \tot\left( \tau_{\mathrm{I}}^{\leq q}X \right) \to \Hm_{\mathrm{I}}^q(X)[-q] \to 0.
	\end{gather*}
	The second is precisely (ii). Part (i) follows from the first together with
	Propositions~\ref{prop:long-exact-sequence-ses} and
	\ref{prop:cone-homotopy}~(i).
\end{proof}

% segment-id: o014.li2.chapter3.double-complex-cohomology.p007
For any object $X$ of $\cate{C}^2(\mathcal{A})$, let
$\Hm_{\mathrm{I}}(X)$ and $\Hm_{\mathrm{II}}(X)$ denote the following double
complexes ($p,q\in\Z$):
\index[sym1]{HmIX@$\Hm_{\mathrm{I}}(X)$, $\Hm_{\mathrm{II}}(X)$}
% segment-id: o014.li2.chapter3.double-complex-cohomology.q009
\begin{equation}\label{eqn:HIHII}\begin{aligned}
	\left( \Hm_{\mathrm{I}}(X)^{p, \bullet}, \dvert_{\Hm_{\mathrm{I}}(X)}^{p, \bullet} \right) & := \Hm_{\mathrm{I}}^p(X), &
	\dhori_{\Hm_{\mathrm{I}}(X)}^{\bullet, \bullet} := 0, \\
	\left( \Hm_{\mathrm{II}}(X)^{\bullet, q}, \dhori_{\Hm_{\mathrm{II}}(X)}^{\bullet, q} \right) & := \Hm_{\mathrm{II}}^q(X), &
	\dvert_{\Hm_{\mathrm{II}}(X)}^{\bullet, \bullet} := 0,
\end{aligned}\end{equation}
In other words, $\Hm_{\mathrm{I}}(X)$ (respectively
$\Hm_{\mathrm{II}}(X)$) is obtained by taking the cohomology of $X$ in the
$\dhori$-direction (respectively the $\dvert$-direction). Both are endofunctors
of $\cate{C}^2(\mathcal{A})$. We may then define
$\Hm_{\mathrm{II}}\Hm_{\mathrm{I}}(X)$ and
$\Hm_{\mathrm{I}}\Hm_{\mathrm{II}}(X)$; both satisfy
$\dhori=0=\dvert$.

% segment-id: o014.li2.chapter3.double-complex-cohomology.th001
\begin{theorem}\label{prop:double-cplx-tot}
	Let $f:X\to Y$ be a morphism in $\cate{C}^2_f(\mathcal{A})$. If the
	induced morphism
	$\Hm_{\mathrm{II}}\Hm_{\mathrm{I}}(X)\to
	\Hm_{\mathrm{II}}\Hm_{\mathrm{I}}(Y)$ is an isomorphism, then
	$\tot(f):\tot(X)\to\tot(Y)$ is a quasi-isomorphism.

	The same conclusion holds if instead the induced morphism
	$\Hm_{\mathrm{I}}\Hm_{\mathrm{II}}(X)\to
	\Hm_{\mathrm{I}}\Hm_{\mathrm{II}}(Y)$ is assumed to be an isomorphism.
\end{theorem}
% segment-id: o014.li2.chapter3.double-complex-cohomology.pf004
\begin{proof}
	The condition is equivalent to saying that the morphism
	$\Hm_{\mathrm{I}}^n(f):\Hm_{\mathrm{I}}^n(X)\to
	\Hm_{\mathrm{I}}^n(Y)$ in $\cate{C}(\mathcal{A})$ is a quasi-isomorphism
	for every $n\in\Z$. The first step is to reduce to the case in which
% segment-id: o014.li2.chapter3.double-complex-cohomology.q010
	\begin{equation}\label{eqn:double-cplx-tot-aux}
		q \ll 0 \implies \tilde{\tau}_{\mathrm{I}}^{\leq q}(X) = 0 = \tilde{\tau}_{\mathrm{I}}^{\leq q}(Y).
	\end{equation}
	Indeed, let $r\in\Z$. Functorial truncation as in
	Lemma~\ref{prop:truncation-ses} shows that the original condition implies
	that
	$\Hm_{\mathrm{I}}^n(\tau_{\mathrm{I}}^{\geq r}f):
	\Hm_{\mathrm{I}}^n(\tau_{\mathrm{I}}^{\geq r}X)\to
	\Hm_{\mathrm{I}}^n(\tau_{\mathrm{I}}^{\geq r}Y)$ is also a
	quasi-isomorphism for every $n$. For a fixed $n$, however, the definition
	of $\cate{C}^2_f(\mathcal{A})$ shows that when $r\ll0$ there is a
	commutative diagram
% segment-id: o014.li2.chapter3.double-complex-cohomology.g003
	\begin{equation}\label{eqn:double-cplx-tot-aux2}\begin{tikzcd}
		\Hm^n(\tot(X)) \arrow[r, "{\Hm^n(\tot(f))}"] \arrow[d, "\sim"' sloped] & \Hm^n(\tot(Y)) \arrow[d, "\sim" sloped] \\
		\Hm^n\left( \tot(\tau_{\mathrm{I}}^{\geq r} X) \right) \arrow[r, "{\Hm^n\left( \tot (\tau_{\mathrm{I}}^{\geq r} f) \right)}"' inner sep=0.8em] & \Hm^n\left( \tot(\tau_{\mathrm{I}}^{\geq r} Y) \right)
	\end{tikzcd}\end{equation}
	The double complexes in the second row satisfy
	\eqref{eqn:double-cplx-tot-aux}. We may therefore replace $f$ by
	$\tau_{\mathrm{I}}^{\geq r}f$ from now on and assume
	\eqref{eqn:double-cplx-tot-aux}. Our goal is to prove that
	$\Hm^n(\tot(f))$ is an isomorphism for the chosen $n$.

	For every $q\in\Z$, Lemma~\ref{prop:double-cplx-tot-aux}~(ii) gives a
	commutative diagram with exact rows
% segment-id: o014.li2.chapter3.double-complex-cohomology.g004
	\[\begin{tikzcd}
		0 \arrow[r] & \tot\left( \tilde{\tau}_{\mathrm{I}}^{\leq q-1} X \right) \arrow[r] \arrow[d, "{\alpha_{q-1}}"] & \tot\left( \tau_{\mathrm{I}}^{\leq q} X \right) \arrow[d, "{\beta_q}"] \arrow[r] & \Hm_{\mathrm{I}}^q(X)[-q] \arrow[d, "{\Hm_{\mathrm{I}}^q(f)[-q]}" description] \arrow[r] & 0 \\
		0 \arrow[r] & \tot\left( \tilde{\tau}_{\mathrm{I}}^{\leq q-1} Y \right) \arrow[r] & \tot\left( \tau_{\mathrm{I}}^{\leq q} Y \right) \arrow[r] & \Hm_{\mathrm{I}}^q(Y)[-q] \arrow[r] & 0
	\end{tikzcd}\]%
	\footnote{Translator's note (O014-C055): the source labels the right-hand
	vertical arrow only by $\simeq$. The theorem's hypothesis says that this
	morphism is the quasi-isomorphism $\Hm_{\mathrm{I}}^q(f)[-q]$, not
	necessarily an isomorphism of complexes; the morphism has therefore been
	named and its precise property stated.}
	The right-hand vertical arrow is a quasi-isomorphism, and all the vertical
	arrows come from $f$. We shall prove that both $\alpha_q$ and $\beta_q$ are
	quasi-isomorphisms. If $\alpha_{q-1}$ is a quasi-isomorphism, then
	functoriality of the associated long exact sequences
	(Proposition~\ref{prop:long-exact-sequence-ses}), together with
	Proposition~\ref{prop:5-lemma}, implies that $\beta_q$ is also a
	quasi-isomorphism. On the other hand,
	Lemma~\ref{prop:double-cplx-tot-aux}~(i) gives the commutative diagram
% segment-id: o014.li2.chapter3.double-complex-cohomology.g005
	\[\begin{tikzcd}
		\tot\left( \tau_{\mathrm{I}}^{\leq q} X \right) \arrow[r, "\text{quasi-isomorphism}"] \arrow[d, "{\beta_q}"'] & \tot\left( \tilde{\tau}_{\mathrm{I}}^{\leq q} X \right) \arrow[d, "{\alpha_q}"] \\
		\tot\left( \tau_{\mathrm{I}}^{\leq q} Y \right) \arrow[r, "\text{quasi-isomorphism}"'] & \tot\left( \tilde{\tau}_{\mathrm{I}}^{\leq q} Y \right)
	\end{tikzcd}\]
	Thus, if $\beta_q$ is a quasi-isomorphism, then so is $\alpha_q$. Since
	$\alpha_q$ is simply $0\rightiso0$ when $q\ll0$, recursion shows that
	$\alpha_q$ and $\beta_q$ are quasi-isomorphisms for all $q$.

	Recall that $n\in\Z$ has been fixed. By an argument similar to that for
	\eqref{eqn:double-cplx-tot-aux2}, when $q\gg0$ there is a commutative
	diagram
% segment-id: o014.li2.chapter3.double-complex-cohomology.g006
	\[\begin{tikzcd}
		\Hm^n(\tot(X)) \arrow[r, "{\Hm^n(\tot(f))}"] & \Hm^n(\tot(Y)) \\
		\Hm^n\left( \tot(\tau_{\mathrm{I}}^{\leq q} X) \right) \arrow[r, "{\Hm^n(\beta_q)}"' inner sep=0.8em, "\sim"] \arrow[u, "\sim" sloped] & \Hm^n\left( \tot(\tau_{\mathrm{I}}^{\leq q} Y) \right) \arrow[u, "\sim"' sloped]
	\end{tikzcd}\]
	Hence $\Hm^n(\tot(f))$ is indeed an isomorphism.

	Finally, the same argument applies when
	$\Hm_{\mathrm{I}}\Hm_{\mathrm{II}}(X)\rightiso
	\Hm_{\mathrm{I}}\Hm_{\mathrm{II}}(Y)$. Alternatively, one may use
	Proposition~\ref{prop:double-cplx-swap} to interchange
	$\Hm_{\mathrm{I}}$ and $\Hm_{\mathrm{II}}$.
\end{proof}

% segment-id: o014.li2.chapter3.double-complex-cohomology.p008
The proof above is rather involved. Example~\sourcecrossref{eg:double-cplx-tot-ss}{in \S~5.6}
will give another proof based on spectral sequences.

% segment-id: o014.li2.chapter3.double-complex-cohomology.co001
\begin{corollary}\label{prop:acyclic-assembly}
	Let $X$ be an object of $\cate{C}_f^2(\mathcal{A})$. If $X$ is row-exact,
	that is, if $\left(X^{\bullet,q},\dhori\right)$ is exact for every
	$q\in\Z$, then $\tot(X)$ is exact. Likewise, if $X$ is column-exact, then
	$\tot(X)$ is exact.
\end{corollary}
% segment-id: o014.li2.chapter3.double-complex-cohomology.pf005
\begin{proof}
	If $X$ is row-exact, then
	$\Hm_{\mathrm{II}}\Hm_{\mathrm{I}}(X)=0$. Theorem~\ref{prop:double-cplx-tot}
	therefore implies that $\tot(X)\to\tot(0)=0$ is a quasi-isomorphism; in
	other words, $\tot(X)$ is exact.

	The argument in the column-exact case is identical. Alternatively, the two
	cases can be deduced from one another by
	Proposition~\ref{prop:double-cplx-swap}.
\end{proof}
